% De Domo Sua --- headnote
% de-domo-sua, 29 September 57 BC, before the College of Pontiffs

\textit{Delivered before the College of Pontiffs on 29
September 57 BC, three weeks after Cicero's return from
exile, on the question whether Clodius's consecration of his
Palatine house --- as a Temple of Liberty --- could stand. The
\textit{religious} half of the question (proper rites,
proper words) Cicero leaves to the pontiffs as their own:
the speech is built almost entirely on \textit{public law}
(\textsection 32). Three threads are then woven: that
Clodius was no lawful tribune at all (so his bill against
Cicero was no law); that even if a lawful tribune, the bill
was a privilegium against an uncondemned citizen (which the
Twelve Tables forbid); and that even if a lawful bill, it
contained no word of consecration --- so the dedication is
the act of one man, not the people, and falls under the
\textit{lex Papiria}, which forbids the consecration of any
building without an order of the plebs.}

\textit{The opening \textsection 1--32 is dressed as
self-defence on a different question --- Cicero's Senate
motion of 7 September which gave Pompey the
\textit{cura annonae}, against which Clodius had heckled in
his preceding speech. This is the answer-to-Clodius preface;
\textsection 32 announces the pivot to the public-law
argument proper. \textsection 34--42 attack Clodius's
\textit{transitio ad plebem} of 58 BC: the adoption was
unlawful (the adopter younger than the adopted, the rite
performed inside three hours and against the auspices), so
Clodius was never tribune. \textsection 43--57 develop the
\textit{privilegium} argument with the famous formulation
that the worst part of the Sullan proscriptions was the
condemnation of citizens by name without trial, and the
mocking treatment of the bill's perverse self-dissolving
drafting (\textsection 47, ``you have been denied'' rather
than ``you are to be denied''). \textsection 58--74 are the
roll-call of those who suffered, the distinction between
the popular Roman people and the rabble of hired men, and
the historical parallels of restored exiles
(Camillus, Servilius Ahala, Popilius, Metellus Numidicus,
Marius). \textsection 95--100 is the most pointed of the
self-defences: Cicero's withdrawal was not flight from
charge, nor fear of the people's vote, nor cowardice ---
but the refusal to drag good men to death along with him.}

\textit{The second half of the speech, \textsection
101--147, is the argument on the house itself. The
historical houses razed for kingship and crime
(\textsection 101: Maelius, Cassius, Vaccus, Manlius);
the satirical genealogy of ``Liberty'' --- the marble image
of a Tanagran courtesan brought back by Q. Catulus's brother
from the looting of Greece (\textsection 111--112); the two
Catuli, father and son, on the overthrow of Catulus's
portico (\textsection 113--114); the murder of Q. Seius for
the adjoining property (\textsection 115); and at the heart
of the legal argument, the parallel of C. Cassius the
censor's referral to the college on the dedication of
Concord, and of the Vestal Licinia's altar struck down
(\textsection 130--137). The closing \textsection 142--147
is the great prayer to the gods of the city --- Jupiter
Capitolinus, Juno Regina, Minerva, the household gods, Mother
Vesta. It worked: the pontiffs would rule that the
consecration was invalid, the Senate would order
reconstruction at public expense, and Cicero would be back
on the Palatine before winter --- though Clodius's gangs
would tear down the rebuilding twice in the following
months.}
